pdfoutput=1
\pdfoutput=1
% ===========================================================================
%  Situs: Physics-Anchored Positional Encoding for State-Conditioned Expertise
%  Target: arXiv → JMLR / NeurIPS Theory Track
%  Date: 2026-06-29
% ===========================================================================

\documentclass[11pt,a4paper]{article}

% ---- Chinese & Font Support ----

% ---- Mathematics ----
\usepackage{amsmath,amssymb,amsthm,mathtools}
\usepackage{bm}
\usepackage{bbm}
\usepackage{dsfont}

% ---- Theorem Environments ----
\newtheorem{theorem}{Theorem}[section]
\newtheorem{proposition}[theorem]{Proposition}
\newtheorem{corollary}[theorem]{Corollary}
\newtheorem{lemma}[theorem]{Lemma}
\newtheorem{definition}[theorem]{Definition}
\newtheorem{assumption}[theorem]{Assumption}
\theoremstyle{remark}
\newtheorem{remark}[theorem]{Remark}
\newtheorem{example}[theorem]{Example}

% ---- Layout ----
\usepackage[margin=2.5cm]{geometry}
\usepackage{booktabs}
\usepackage{graphicx}
\usepackage{hyperref}
\usepackage{enumitem}
\usepackage[capitalise,noabbrev]{cleveref}

% ---- Math Operators ----
\DeclareMathOperator{\PE}{PE}
\DeclareMathOperator{\Cov}{Cov}
\DeclareMathOperator{\Var}{Var}
\DeclareMathOperator{\Corr}{Corr}
\DeclareMathOperator*{\argmin}{arg\,min}
\DeclareMathOperator*{\argmax}{arg\,max}
\DeclareMathOperator{\KL}{D_{KL}}
\DeclareMathOperator{\TV}{TV}

% ---- Custom notations ----
\newcommand{\R}{\mathbb{R}}
\newcommand{\N}{\mathbb{N}}
\newcommand{\Y}{\mathcal{Y}}
\newcommand{\X}{\mathcal{X}}
\newcommand{\Sstates}{\mathcal{S}}
\newcommand{\Ppos}{\mathcal{P}}
\newcommand{\E}{\mathbb{E}}
\newcommand{\V}{\mathbb{V}}
\newcommand{\I}{\mathbb{I}}
\newcommand{\indep}{\perp\!\!\!\perp}
\newcommand{\Situs}{\textsc{Situs}}
\newcommand{\pe}{\mathrm{pe}}

% ---- Honest critique markers ----
\newcommand{\rigorous}{\textnormal{\textbf{[]}}}
\newcommand{\heuristic}{\textnormal{\textbf{[]}}}
\newcommand{\openproblem}{\textnormal{\textbf{[]}}}

\title{\textbf{\LARGE Situs: Physics-Anchored Positional Encoding\\for State-Conditioned Expertise}}

\author{SCX}

\date{\today}

\begin{document}

\maketitle

% ===========================================================================
% ABSTRACT
% ===========================================================================
\begin{abstract}
Positional Encoding, PE——
token\Situs{}——
\textbf{}
\Situs{}
$I(Y; P \mid S) > 0$

\Situs{}\textbf{}Bochner
Laplace1.2.1$O(1/d)$
1.2.23DLipschitz
1.2.31.3.1Lipschitz1.4.1
\textbf{}Theorem~1$\delta_s^{\PE}$
$\Delta_s^{\Situs} = \Delta_s + \delta_s^{\PE}$
$\delta_s^{\PE} > 0$$I(Y; P \mid S) > 0$2.2.1
2.3.1$\delta_s^{\PE}$
Fano2.4.1\textbf{}Theorem~2
$\varepsilon_{\PE} = I(Y; P \mid X) - I(Y; \PE(P) \mid X)$3.1
Theorem~$2'$KSG
$\varepsilon_{\PE}$3.3.13.3.1
\textbf{}Theorem~$3'$-
34.1
\textbf{Open Problem}``''
Theorem~$3'$


3D\Situs{}



\textbf{}Lipschitz
Bochner
\end{abstract}

% ===========================================================================
% SECTION 1: INTRODUCTION
% ===========================================================================
\section{}

\subsection{}

Positional Encoding, PE
Transformer——\textbf{}
\textbf{}token
\cite{vaswani2017attention}

——\cite{vaswani2017attention}
RoPE\cite{su2024roformer}\cite{devlin2019bert}——
\textbf{}

\begin{enumerate}[label=(\roman*)]
    \item 
          ``token''token
    \item ``''\textbf{}——
          
    \item perplexity
          
\end{enumerate}

\textbf{}——
-——
\textbf{}token


\textbf{}Lys
tokenizationtoken
Lys-37$\text{p}K_a$6--8
Lys-289$\text{p}K_a \approx 10.4$
\textbf{}
Lys\textbf{}


\subsection{\Situs{}}

\Situs{}\textit{situs}``''``''——
\textbf{}\Situs{}

\begin{quote}
    \textbf{\Situs{}}$P$$S$$Y$
    ——$I(Y; P \mid S) > 0$——
    $\PE: \Ppos \to \R^{d_{\pe}}$
    
\end{quote}

\Situs{}1.1.1--1.1.3
\begin{itemize}
    \item \textbf{} $\Ppos_{\text{seq}} = \{1, 2, \ldots, L\}$
          N→C
    \item \textbf{} $\Ppos_{\text{3D}} = \R^3$
          ——
          \textbf{}
    \item \textbf{} $\Ppos_{\text{count}} = \N$
          
\end{itemize}

\subsection{SCXCC}

\Situs{}SCXState-Conditioned eXpertise
SCX\cite{scx2026theorems}$M$
$\{E_1, \ldots, E_M\}$$s_i$$y_i$$m$
$v_m(s_i) = \mathbf{1}[E_m(s_i) \neq y_i] \in \{0, 1\}$
$v_m \sim \text{Bernoulli}(p_{\text{clean}, s})$$p_{\text{clean}, s} < 0.5$
$v_m \sim \text{Bernoulli}(p_{\text{noisy}, s})$$p_{\text{noisy}, s} > 0.5$
$\Delta_s = p_{\text{noisy}, s} - p_{\text{clean}, s} > 0$

SCX\cite{scx2026cc_audit}
\begin{itemize}
    \item \textbf{Theorem 1}$F_1 \geq 1 - \frac{1}{\eta}\sum_s \rho_s \exp(-2M\Delta_s^2)$Chernoff-Hoeffding
    \item \textbf{Theorem 2}$F_{1,\text{SCX}} \leq F_{1,\text{base}} + C_F\sqrt{\delta/2}$Fano
    \item \textbf{Theorem 3}$W_A, W_B$
    \item \textbf{Theorem 4}MinimaxChernoff-Stein
\end{itemize}

CCCorrectness \& Completeness\cite{scx2026cc_audit}
Spring
——sign error

\textbf{}
\begin{enumerate}[label=(\arabic*)]
    \item \textbf{}\S2
          Lipschitz
    \item \textbf{Theorem 1}\S3CC
          $\delta_s^{\PE} > 0$DPIFano
    \item \textbf{Theorem 2}\S4$\varepsilon_{\PE}$
          
    \item \textbf{Theorem 3$'$}\S5PEPE
          PEPE
          ``''
    \item \textbf{Honest Critique}
          IDR
\end{enumerate}

% ===========================================================================
% SECTION 2: ENCODING FUNCTIONS
% ===========================================================================
\section{}

\Situs{}
\textbf{}
Lipschitz

% ---------------------------------------------------------------------------
\subsection{}

\begin{definition}[——]
\label{def:1.1.1}
$L \in \N$
\begin{equation}
    \Ppos_{\text{seq}} = \{1, 2, \ldots, L\} \subset \N
\end{equation}
$\Ppos_{\text{seq}} \cong [0, L] \subset \R$\\
\textbf{}NC
\end{definition}

\begin{definition}[——]
\label{def:1.1.2}

\begin{equation}
    \Ppos_{\text{3D}} = \R^3
\end{equation}
$\mathbf{p} = (x, y, z) \in \R^3$\\
\textbf{}
\textbf{}
\end{definition}

\begin{definition}[——]
\label{def:1.1.3}

\begin{equation}
    \Ppos_{\text{count}} = \N
\end{equation}
$N$\\
\textbf{}
\end{definition}

% ---------------------------------------------------------------------------
\subsection{}

\subsubsection{}

\begin{definition}[]
\label{def:1.2.1}
$p \in [0, L]$$d$
\begin{equation}
    \boxed{
    \PE_{\text{scalar}}(p, 2j) = \sqrt{\frac{2}{d}}\,\sin\!\left(\frac{2\pi p}{\lambda_j}\right),\quad
    \PE_{\text{scalar}}(p, 2j+1) = \sqrt{\frac{2}{d}}\,\cos\!\left(\frac{2\pi p}{\lambda_j}\right)
    }
\end{equation}
$j = 0, 1, \ldots, d/2-1$$\lambda_j > 0$\textbf{}\\
$\sqrt{2/d}$$\|\PE_{\text{scalar}}(p)\| = 1$$p$
$d/2$$(\sin,\cos)$$(2/d) \cdot 1 = 2/d$$1$
$d \to \infty$\textbf{}$k(\Delta)$$O(d)$\\
\textbf{}Transformer$\omega_j = 10000^{-2j/d}$
\textbf{}$\lambda_j$
$\lambda_j$
\end{definition}


\begin{equation}
    \langle \PE_{\text{scalar}}(p), \PE_{\text{scalar}}(q) \rangle
    = \frac{2}{d}\sum_{j=0}^{d/2-1} \cos\!\left(\frac{2\pi(p-q)}{\lambda_j}\right)
    = \frac{2}{d}K_{\PE}(\Delta)
\end{equation}
$\Delta = p - q$
$\Delta = 0$$\langle \PE(p), \PE(p) \rangle = (2/d) \cdot (d/2) = 1 = k(0)$
$d \to \infty$$(2/d)K_{\PE}(\Delta) \to \int_0^\infty S(\omega)\cos(\omega\Delta)d\omega = k(\Delta)$

\subsubsection{}

$d$$L$$\{\lambda_j\}_{j=0}^{d/2-1}$
``''

\begin{assumption}[]
\label{ass:1.2.1}
\textbf{}
$k(\Delta)$$\Delta = |p-q|$
\begin{equation}
    k(\Delta) = \exp\!\left(-\frac{\Delta}{\xi}\right)
\end{equation}
$\xi > 0$\textbf{}$\approx 1.9$\,\AA{}
BochnerFourier
\end{assumption}

Laplace
FriedelCoulomb
Yukawa$\xi$Thomas-Fermi

\begin{theorem}[——]
\label{thm:1.2.1}
$k(\Delta) = \exp(-|\Delta|/\xi)$Laplace
$\frac{2}{d}K_{\PE}(\Delta) = \frac{2}{d}\sum_{j=0}^{d/2-1} \cos(2\pi\Delta/\lambda_j)$
$L^2([0, L])$$k(\Delta)$$\{\lambda_j\}$
\begin{equation}
    \boxed{
    \lambda_j = 2\pi\xi \cdot \cot\!\left(\frac{\pi(2j+1)}{2d}\right),\quad
    j = 0, 1, \ldots, \frac{d}{2}-1
    }
\end{equation}
\end{theorem}

\begin{proof}
\textbf{1Bochner}LaplaceFourier
\begin{equation}
    \hat{k}(\omega) = \int_{-\infty}^{\infty} e^{-i\omega\Delta} e^{-|\Delta|/\xi}\,d\Delta
    = \frac{2\xi}{1 + \xi^2\omega^2}
\end{equation}
Cauchy$1/\xi$

\textbf{2}Fourier$\hat{k}$
\begin{equation}
    k(\Delta) = \frac{1}{2\pi}\int_{-\infty}^{\infty} e^{i\omega\Delta}\hat{k}(\omega)\,d\omega
    = \int_{0}^{\infty} \cos(\omega\Delta) \cdot S(\omega)\,d\omega
\end{equation}
$S(\omega) = \frac{2\xi}{\pi(1 + \xi^2\omega^2)}$
$\int_0^{\infty} S(\omega)\,d\omega = 1$

\textbf{3}$\frac{2}{d}K_{\PE}(\Delta) = \frac{2}{d}\sum_{j=0}^{d/2-1} \cos(\omega_j \Delta)$
$\omega_j = 2\pi/\lambda_j$\textbf{}$2/d$
$\int_0^{\infty} S(\omega)\cos(\omega\Delta)\,d\omega$
$\{S(\omega_j)\Delta\omega_j\}$
$Q(\omega) = \int_0^{\omega} S(u)\,du = \frac{2}{\pi}\arctan(\xi\omega)$
\textbf{}

\textbf{4}$j$$(2j+1)/d$
\begin{equation}
    Q(\omega_j) = \frac{2j+1}{d} \quad\Longrightarrow\quad
    \omega_j = \frac{1}{\xi}\tan\!\left(\frac{\pi(2j+1)}{2d}\right)
\end{equation}
$\lambda_j = 2\pi/\omega_j$\hfill $\square$
\end{proof}

\rigorous{} \textbf{}Gauss quadrature
——BochnerFourier——


\begin{corollary}[]
\label{cor:1.2.1}
$j=0$
\begin{equation}
    \lambda_{\max} = 2\pi\xi \cdot \cot\!\left(\frac{\pi}{2d}\right)
    \approx 4d\xi \quad (\text{ } d)
\end{equation}
$\lambda_{\max} \geq 2L$Nyquist
\begin{equation}
    \boxed{d_{\min} \approx \frac{L}{2\xi}}
\end{equation}
\end{corollary}

\begin{corollary}[Transformer]
\label{cor:1.2.2}
Transformer$\omega_j \propto 10000^{-2j/d}$
\textbf{}Laplace
$\omega_j \propto \tan(\pi(2j+1)/2d)$\textbf{}
NLP
\end{corollary}

\begin{theorem}[]
\label{thm:1.2.2}
Laplace$d/2$
\begin{equation}
    \boxed{
    \sup_{\Delta \in [0, L]} \left|\frac{2}{d}K_{\PE}(\Delta) - k(\Delta)\right|
    = O\!\left(\frac{1}{d}\right)
    }
\end{equation}
$d \to \infty$$O(1/d)$$\xi$$L$$O(1/d)$
$S(\omega)$$k(\Delta)$$[0,L]$
\end{theorem}

\rigorous{} \textbf{}
Fourier
Fourier
$(2/d)K_{\PE}$$K_{\PE}$$k(\Delta)$
$O(1)$$K_{\PE}(0)=d/2$$k(0)=1$

% ---------------------------------------------------------------------------
\subsection{Lipschitz}

\begin{theorem}[SitusLipschitz ()]
\label{thm:1.2.3}
$\PE_{\text{scalar}}: [0, L] \to \R^d$Lipschitz
\begin{equation}
    \boxed{
    L_{\PE}^{\text{scalar}} = \frac{2\sqrt{2}\,\pi}{\sqrt{d}} \cdot \sqrt{\sum_{j=0}^{d/2-1}\frac{1}{\lambda_j^2}}
    }
\end{equation}
\end{theorem}

\begin{proof}
$(2j, 2j+1)$$\PE_j(p) = \sqrt{2/d}\,(\sin(2\pi p/\lambda_j), \cos(2\pi p/\lambda_j))$
\begin{align}
    \|\PE_j(p) - \PE_j(q)\|^2
    &= \frac{2}{d}\Bigl[\bigl[\sin(2\pi p/\lambda_j) - \sin(2\pi q/\lambda_j)\bigr]^2
     + \bigl[\cos(2\pi p/\lambda_j) - \cos(2\pi q/\lambda_j)\bigr]^2\Bigr] \nonumber\\
    &= \frac{2}{d} \cdot 4\sin^2\!\left(\frac{\pi(p-q)}{\lambda_j}\right) \nonumber\\
    &\leq \frac{2}{d} \cdot \left(\frac{2\pi|p-q|}{\lambda_j}\right)^2
     = \frac{8\pi^2|p-q|^2}{d \cdot \lambda_j^2}
\end{align}
$(\sin a - \sin b)^2 + (\cos a - \cos b)^2 = 4\sin^2((a-b)/2)$
$|\sin\theta| \leq |\theta|$
$j$$p, q$$\sin\theta \approx \theta$\hfill $\square$
\end{proof}

\rigorous{} \textbf{}$p, q$


\begin{corollary}[Lipschitz]
\label{cor:1.2.3}

\begin{equation}
    L_{\PE}^{\text{scalar}} \sim \frac{1}{\xi} \cdot \sqrt{\frac{d}{6}} \quad (d \to \infty)
\end{equation}
Lipschitz$O(d)$$O(\sqrt{d})$——$\sqrt{2/d}$
$\sqrt{d}$\textbf{}
\end{corollary}

% ---------------------------------------------------------------------------
\subsection{}

\subsubsection{}

\begin{definition}[]
\label{def:1.3.1}
$\mathbf{p} = (x, y, z) \in \R^3$$d$$d \geq 6$
\begin{equation}
    \boxed{
    \PE_{\text{rot}}(\mathbf{p}) = \mathbf{R}(\mathbf{p}) \cdot \mathbf{e}_0
    }
\end{equation}
$\mathbf{e}_0 \in \R^d$
$\mathbf{R}(\mathbf{p}) = \mathbf{R}_x(\alpha x) \cdot \mathbf{R}_y(\beta y) \cdot \mathbf{R}_z(\gamma z) \in SO(d)$
$\alpha, \beta, \gamma > 0$\textbf{}rad/\AA{}
$\mathbf{R}_a(\theta)$$d$
$(0,1)$, $(2,3)$, $(4,5)$$d-6$
\end{definition}

\subsubsection{}

\begin{proposition}[$\R^3 \to SO(d)$]
\label{prop:1.3.1}
$\mathbf{p} \mapsto \mathbf{R}(\mathbf{p})$$\R^3$$SO(d)$Abel
\textbf{}
\end{proposition}


\begin{equation}
    [\mathbf{R}_x(\theta), \mathbf{R}_y(\phi)] = 0,\quad
    [\mathbf{R}_y(\phi), \mathbf{R}_z(\psi)] = 0,\quad
    [\mathbf{R}_z(\psi), \mathbf{R}_x(\theta)] = 0
\end{equation}
$\mathbf{R}(\mathbf{p}_1 + \mathbf{p}_2) = \mathbf{R}(\mathbf{p}_1) \cdot \mathbf{R}(\mathbf{p}_2)$

\rigorous{} \textbf{}

\subsubsection{}

$\mathbf{p}, \mathbf{q} \in \R^3$
\begin{align}
    \langle \PE_{\text{rot}}(\mathbf{p}), \PE_{\text{rot}}(\mathbf{q}) \rangle
    &= \langle \mathbf{R}(\mathbf{p})\mathbf{e}_0, \mathbf{R}(\mathbf{q})\mathbf{e}_0 \rangle
     = \langle \mathbf{e}_0, \mathbf{R}(\mathbf{q} - \mathbf{p})\mathbf{e}_0 \rangle \nonumber\\
    &= \cos(\alpha\Delta x) + \cos(\beta\Delta y) + \cos(\gamma\Delta z)
\end{align}
$\Delta\mathbf{p} = \mathbf{q} - \mathbf{p}$\textbf{}
——


\subsubsection{Lipschitz}

\begin{theorem}[Lipschitz ()]
\label{thm:1.3.1}
$\PE_{\text{rot}}: \R^3 \to \R^d$Lipschitz$\mathbf{e}_0$
1
\begin{equation}
    \boxed{
    L_{\PE}^{\text{rot}} = \max(\alpha, \beta, \gamma)
    }
\end{equation}
\end{theorem}

\begin{proof}
\textbf{1}
$\PE_{\text{rot}}(\mathbf{p}) - \PE_{\text{rot}}(\mathbf{q})
 = \mathbf{R}(\mathbf{p})\mathbf{e}_0 - \mathbf{R}(\mathbf{q})\mathbf{e}_0$

\begin{align}
    \|\PE_{\text{rot}}(\mathbf{p}) - \PE_{\text{rot}}(\mathbf{q})\|^2
    &= 2 - 2\cos(\alpha\Delta x) + 2 - 2\cos(\beta\Delta y) + 2 - 2\cos(\gamma\Delta z) \nonumber\\
    &= 4\sin^2\!\left(\frac{\alpha\Delta x}{2}\right)
     + 4\sin^2\!\left(\frac{\beta\Delta y}{2}\right)
     + 4\sin^2\!\left(\frac{\gamma\Delta z}{2}\right)
\end{align}

\textbf{2$|\sin\theta| \leq |\theta|$}
\begin{equation}
    \|\PE_{\text{rot}}(\mathbf{p}) - \PE_{\text{rot}}(\mathbf{q})\|^2
    \leq \alpha^2\Delta x^2 + \beta^2\Delta y^2 + \gamma^2\Delta z^2
    \leq \max(\alpha, \beta, \gamma)^2 \cdot \|\Delta\mathbf{p}\|_2^2
\end{equation}
$\|\PE_{\text{rot}}(\mathbf{p}) - \PE_{\text{rot}}(\mathbf{q})\|
 \leq \max(\alpha, \beta, \gamma) \cdot \|\mathbf{p} - \mathbf{q}\|_2$

\textbf{3}$\Delta\mathbf{p} = (\varepsilon, 0, 0)$$\alpha = \max(\alpha, \beta, \gamma)$
$\varepsilon \to 0$$\|\PE_{\text{rot}}(\mathbf{p}) - \PE_{\text{rot}}(\mathbf{q})\|
= 2|\sin(\alpha\varepsilon/2)| \to \alpha|\varepsilon|
= \max(\alpha, \beta, \gamma) \cdot \|\Delta\mathbf{p}\|$\hfill $\square$
\end{proof}

\rigorous{} \textbf{}
$\sqrt{3}$

\begin{remark}[]
\label{rem:1.3.1}
$\|\mathbf{R}(\mathbf{p}) - \mathbf{R}(\mathbf{q})\|_F$
Cauchy-Schwarz$2\sqrt{\alpha^2 + \beta^2 + \gamma^2}$
$\max(\alpha, \beta, \gamma)$$2\sqrt{3} \approx 3.46$$\alpha=\beta=\gamma$
i$\|D_x+D_y+D_z\|_F
\leq \|D_x\|_F + \|D_y\|_F + \|D_z\|_F$$\sqrt{3}$
iiCauchy-Schwarz$\ell_1$$\ell_2$$\sqrt{3}$
$\ell_2$
\end{remark}

% ---------------------------------------------------------------------------
\subsection{Lipschitz}

\begin{theorem}[Lipschitz ()]
\label{thm:1.4.1}
$\PE: \Ppos \to \R^{d_{\pe}}$$L_{\PE} > 0$
$\mathbf{p}, \mathbf{q} \in \Ppos$
\begin{equation}
    \boxed{
    \|\PE(\mathbf{p}) - \PE(\mathbf{q})\| \leq L_{\PE} \cdot d_{\Ppos}(\mathbf{p}, \mathbf{q})
    }
\end{equation}
$d_{\Ppos}$$\Ppos$$|p-q|$3D$\|\mathbf{p} - \mathbf{q}\|_2$
~\ref{tab:lipschitz}
\end{theorem}

\begin{table}[htbp]
\centering
\caption{Lipschitz}
\label{tab:lipschitz}
\begin{tabular}{@{}lll@{}}
\toprule
\textbf{} & $\Ppos$ & $L_{\PE}$ \\
\midrule
 & $[0, L]$ & $\frac{2\sqrt{2}\,\pi}{\sqrt{d}}\sqrt{\sum_j 1/\lambda_j^2}$ \\
3D & $\R^3$ & $\max(\alpha, \beta, \gamma)$ \\
\bottomrule
\end{tabular}
\end{table}

\textbf{}$L_{\PE}$\textbf{}$L_{\PE}$
$L_{\PE}$
$\xi$

% ===========================================================================
% SECTION 3: THEOREM 1 CORRECTION
% ===========================================================================
\section{Theorem 1}

% ---------------------------------------------------------------------------
\subsection{——}

\Situs{}$m$$s_i$
\begin{equation}
    h_i = \phi(s_i) + \PE(p_i)
\end{equation}
$\phi: \R^{d_s} \to \R^d$

\begin{assumption}[]
\label{ass:2.1.1}
$E_m$$\R^d$$h_i$
$v_m^{\Situs}(s_i) = \mathbf{1}[E_m(h_i) \neq y_i]$
/$p_{\text{clean}, s}^{\Situs}$$p_{\text{noisy}, s}^{\Situs}$
\end{assumption}

\begin{proposition}[\Situs{}——CC]
\label{prop:2.1}
\Situs{}
\begin{equation}
    \boxed{
    \Delta_s^{\Situs} = \Delta_s + \delta_s^{\PE}
    }
\end{equation}

\begin{equation}
    \boxed{
    \delta_s^{\PE} = \underbrace{(p_{\text{noisy}, s}^{\Situs} - p_{\text{noisy}, s})}_{
    \text{}}
    - \underbrace{(p_{\text{clean}, s}^{\Situs} - p_{\text{clean}, s})}_{
    \text{}}
    }
\end{equation}
$\delta_s^{\PE}$$[-\Delta_s, 1 - p_{\text{clean}, s}]$
\end{proposition}

\begin{proof}

\begin{align}
    \Delta_s^{\Situs} &= p_{\text{noisy}, s}^{\Situs} - p_{\text{clean}, s}^{\Situs} \nonumber\\
    &= (p_{\text{noisy}, s} + (p_{\text{noisy}, s}^{\Situs} - p_{\text{noisy}, s}))
     - (p_{\text{clean}, s} + (p_{\text{clean}, s}^{\Situs} - p_{\text{clean}, s})) \nonumber\\
    &= \Delta_s + \delta_s^{\PE}
\end{align}
$p_{\text{noisy}}^{\Situs}, p_{\text{clean}}^{\Situs} \in [0,1]$\hfill $\square$
\end{proof}

\rigorous{} \textbf{}CC\S2.32.1$\delta_s^{\PE}$
\textbf{}$\delta_s^{\PE} > 0$\Situs{}


\textbf{Theorem 1\Situs{}}
\begin{equation}
    \boxed{
    F_1^{\Situs} \geq 1 - \frac{1}{\eta} \sum_{s \in \Sstates} \rho_s \cdot
    \exp\!\left(-2M (\Delta_s + \delta_s^{\PE})^2\right)
    }
\end{equation}
Chernoff-Hoeffding$[0,1]$$\Delta_s$
$\Delta_s + \delta_s^{\PE}$

% ---------------------------------------------------------------------------
\subsection{$\delta_s^{\PE} > 0$}

\begin{theorem}[$\delta_s^{\PE} > 0$——]
\label{thm:2.2.1}
$S, P, Y$
\begin{equation}
    \boxed{
    I(Y; P \mid S) > 0
    }
\end{equation}
\textbf{}\Situs{}
$\PE: \Ppos \to \R^{d_{\pe}}$$d_{\pe}$$s$
\textbf{}$(s, \PE(p))$$\delta_s^{\PE} > 0$
\end{theorem}

\begin{proof}
\textbf{1}$I(Y; P \mid S) > 0$
$\mathcal{A} \subset \Sstates \times \Ppos$$P_{Y|S=s, P=p} \neq P_{Y|S=s}$

\textbf{2}$\PE$
\begin{equation}
    I(Y; \PE(P) \mid S) = I(Y; P \mid S) - \varepsilon_{\PE}
\end{equation}
$\varepsilon_{\PE} \geq 0$\S4$\PE$
$\varepsilon_{\PE} < I(Y; P \mid S)$$I(Y; \PE(P) \mid S) > 0$

\textbf{3}
$y$$(s, p)$$y$
→$p_{\text{noisy}, s}^{\Situs} > p_{\text{noisy}, s}$
→
$p_{\text{clean}, s}^{\Situs} \leq p_{\text{clean}, s}$
\textbf{}$\delta_s^{\PE} > 0$

\textbf{4——}\heuristic{} 
SCX$E_m$$\delta_s^{\PE} > 0$
$\ell_1$$P_{Y|S,P}$
$|\delta_s^{\PE}|$\textbf{}
\textbf{}1--3
\heuristic{}\hfill $\square$
\end{proof}

\textbf{}1--34
\heuristic{}——
VCRademacher
42.2.1$I(Y;P|S)>0$
$P$——

\begin{proposition}[$\delta_s^{\PE} > 0$]
\label{prop:2.2.1}
$m$$\delta_s^{\PE} > 0$
\begin{enumerate}[label=(\roman*)]
    \item $I(Y; P \mid S) > 0$——
    \item \Situs{}-——
\end{enumerate}
$(i)$
\begin{equation}
    \exists s, p_1, p_2 \text{  } P_{Y|S=s, P=p_1} \neq P_{Y|S=s, P=p_2}
\end{equation}
\end{proposition}

\textbf{}
\begin{itemize}
    \item $s$$\alpha$-$p_1$$\beta$-
          $p_2$$\implies I(Y; P \mid S) > 0$
    \item $s$$p_1$$p_2$
          $\implies I(Y; P \mid S) > 0$
\end{itemize}

% ---------------------------------------------------------------------------
\subsection{$\delta_s^{\PE}$}

\begin{theorem}[$\delta_s^{\PE}$]
\label{thm:2.3.1}
DPIPinsker
\begin{equation}
    \boxed{
    \delta_s^{\PE} \leq \min\!\left(1 - p_{\text{clean}, s},\,
    \sqrt{\frac{1}{2} D_{\KL}(P_{\PE|S,Y} \| P_{\PE|S})}\right)
    + \delta_s^{\text{variance}}
    }
\end{equation}
$\delta_s^{\text{variance}} = O(1/\sqrt{M})$$M$
$I(P; Y \mid S) = 0$$M \to \infty$$\delta_s^{\PE} \leq 0$
\end{theorem}

\begin{proof}
\textbf{1Pinsker}$A$$|P(A) - Q(A)| \leq \sqrt{\frac{1}{2}D_{\KL}(P\|Q)}$

\begin{align}
    |p_{\text{clean}, s}^{\Situs} - p_{\text{clean}, s}|
    &\leq \sqrt{\frac{1}{2} D_{\KL}(P_{\PE(P)|S=s, \text{clean}} \| P_{\PE(P)|S=s})}\\
    |p_{\text{noisy}, s}^{\Situs} - p_{\text{noisy}, s}|
    &\leq \sqrt{\frac{1}{2} D_{\KL}(P_{\PE(P)|S=s, \text{noisy}} \| P_{\PE(P)|S=s})}
\end{align}

\textbf{2}
$\delta_s^{\PE} \leq |p_{\text{noisy}, s}^{\Situs} - p_{\text{noisy}, s}|
                      + |p_{\text{clean}, s}^{\Situs} - p_{\text{clean}, s}|$

\textbf{3}$\delta_s^{\PE}$2.1
$\delta_s^{\PE} \leq 1 - p_{\text{clean}, s}$$p_{\text{noisy}}^{\Situs}=1, p_{\text{clean}}^{\Situs}=0$
$[0,1]$PinskerKL
$\min$\textbf{}
$\delta_s^{\PE} \leq 1-p_{\text{clean},s}$
$\delta_s^{\PE} \leq \sqrt{\frac{1}{2}D_{\KL}^{(1)}} + \sqrt{\frac{1}{2}D_{\KL}^{(2)}}$Pinsker
$\delta_s^{\PE} \leq \min(\text{}, \text{Pinsker})$
KLPinsker$p_{\text{clean},s}$

\textbf{4DPI}$I(\PE(P); Y \mid S) \leq I(P; Y \mid S)$
$I(P; Y \mid S) = 0$KL$M \to \infty$$\delta_s^{\PE} \leq 0$
\hfill $\square$
\end{proof}

\rigorous{} \textbf{1--3Pinsker++4
$M \to \infty$——}
$M$$\delta_s^{\text{variance}}$$O(1/\sqrt{M})$
\textbf{}$1-p_{\text{clean},s}$$\delta_s^{\PE}$
2.1Pinsker$\min$\textbf{}
$\min$——$\min$

% ---------------------------------------------------------------------------
\subsection{$\delta_s^{\PE}$Fano}

\begin{proposition}[$\delta_s^{\PE}$——Fano]
\label{prop:2.4.1}
\heuristic{} \textbf{}$P$$S$$Y$
PE$\rho = I(Y; \PE(P) \mid S) / I(Y; P \mid S) \in (0, 1]$
Fano
$s$
\begin{equation}
    \delta_s^{\PE} \approx \frac{2\rho \cdot I(Y; P \mid S) - \log 2}{\log |\Y|}
    - O\!\left(\frac{1}{\sqrt{M}}\right)
\end{equation}
\textbf{}\textbf{}Fano
$P_e \geq A$$P_e^{\Situs} \geq B$$P_e - P_e^{\Situs} \geq A - B$
$a \geq A$$b \geq B$$a-b \geq A-B$Fano
$\delta_s^{\PE}$\textbf{}

\end{proposition}

\begin{proof}[]
\textbf{1Fano}
\begin{equation}
    H(Y \mid S, \PE(P)) \leq H_2(P_e) + P_e \cdot \log(|\Y| - 1)
\end{equation}
$P_e \geq (H(Y \mid S, \PE(P)) - 1) / \log |\Y|$

\textbf{2}Fano

\begin{align}
    P_e - P_e^{\Situs}
    &\stackrel{\text{}}{=} \frac{H(Y|S) - H(Y|S, \PE(P))}{\log |\Y|} \nonumber\\
    &= \frac{\rho \cdot I(Y; P \mid S)}{\log |\Y|}
\end{align}
$\Delta P_e$$\delta_s^{\PE}$2
\textbf{}\hfill $\square$
\end{proof}

\textbf{}1Fano2--3\heuristic{}
(i) $\Delta P_e$$\delta_s^{\PE}$\Situs{}
(ii) Fano——
(iii) $\log 2$nat$\ln 2$ natsbit$1$
FanoFanoFano

\textbf{}$I(Y; P \mid S)$$\delta_s^{\PE}$
——\Situs{}\textbf{
\Situs{}}

% ---------------------------------------------------------------------------
\subsection{}

\begin{theorem}[$\delta_s^{\PE}$]
\label{thm:2.5.1}
PELipschitz$L_{\PE}$$E_m$\textbf{}
$E_m^{\text{soft}}: \R^d \to [0,1]$softmax$L_E$-Lipschitz
\textbf{}$E_m$$h_i, h_j$$\tau > 0$
$|E_m^{\text{soft}}(h) - 0.5| \geq \tau$$h \in \{h_i, h_j\}$
$s$$p_i, p_j$
\begin{equation}
    \boxed{
    |\delta_i^{\PE} - \delta_j^{\PE}|
    \leq 2 \cdot L_E \cdot L_{\PE} \cdot d_{\Ppos}(p_i, p_j)
    + O\!\left(\frac{1}{\sqrt{M}}\right)
    }
\end{equation}
\textbf{}$h_i$$h_j$
non-vacuous——$\mathbf{1}[E_m(h) \neq y]$
$L_E$-Lipschitz$L_E L_{\PE} d_{\Ppos}(p_i, p_j)$
\end{theorem}

\begin{proof}
\begin{align}
    |\delta_i^{\PE} - \delta_j^{\PE}|
    &= |(p_{\text{noisy}, i}^{\Situs} - p_{\text{noisy}, j}^{\Situs})
       - (p_{\text{clean}, i}^{\Situs} - p_{\text{clean}, j}^{\Situs})| \nonumber\\
    &\leq |p_{\text{noisy}, i}^{\Situs} - p_{\text{noisy}, j}^{\Situs}|
       + |p_{\text{clean}, i}^{\Situs} - p_{\text{clean}, j}^{\Situs}| \nonumber\\
    &\leq 2L_E \cdot \|h_i - h_j\|
     = 2L_E \cdot \|\PE(p_i) - \PE(p_j)\|
     \leq 2L_E L_{\PE} \cdot d_{\Ppos}(p_i, p_j)
\end{align}
$s_i = s_j = s$$\phi(s_i) = \phi(s_j)$
\textbf{}$p_{\text{clean},i}^{\Situs}
= \E[\mathbf{1}[E_m(h_i) \neq y]]$\textbf{}
$E_m^{\text{soft}}$$L_E$-Lipschitz$\mathbf{1}[\cdot \neq y]$
$|E_m^{\text{soft}}(h) - 0.5| \geq \tau$$h_i$$h_j$
Lipschitz
$p_{\text{clean},i}^{\Situs}$$p_{\text{clean},j}^{\Situs}$
$\|h_i - h_j\|$
$O(1/\sqrt{M})$\hfill $\square$
\end{proof}

\rigorous{} \textbf{}
vacuousCC2.2


\textbf{}\textbf{}——
\textbf{}
$\delta^{\PE}$(i) 
 (ii) $L_E L_{\PE}$

% ===========================================================================
% SECTION 4: THEOREM 2 CORRECTION
% ===========================================================================
\section{Theorem 2}

% ---------------------------------------------------------------------------
\subsection{$\varepsilon_{\PE}$}


\begin{itemize}
    \item $X$$S$$\phi(S)$
    \item $P$
    \item $\PE(P)$
    \item $Y$
\end{itemize}
Markov
\begin{equation}
    Y \leftrightarrow (X, P) \leftrightarrow (X, \PE(P))
\end{equation}
$Y \indep \PE(P) \mid (X, P)$——

\begin{definition}[$\varepsilon_{\PE}$——]
\label{def:3.1}
\begin{equation}
    \boxed{
    \varepsilon_{\PE} = I(Y; P \mid X) - I(Y; \PE(P) \mid X)
    }
\end{equation}
\end{definition}

\textbf{}
\begin{enumerate}[label=(\arabic*)]
    \item \textbf{DPI}$\varepsilon_{\PE} \geq 0$
          $(X, \PE(P))$$(X, P)$
    \item \textbf{}$\varepsilon_{\PE} = 0$PE$X$$Y$
          \textbf{}$Y \indep P \mid (X, \PE(P))$
    \item \textbf{}$I(Y; P \mid X) = 0$$\varepsilon_{\PE} = 0$
    \item \textbf{}$\varepsilon_{\PE} \leq I(Y; P \mid X) \leq H(Y)$
\end{enumerate}

\rigorous{} \textbf{}(1)DPI
(2)(3)(4)

\begin{proposition}[-]
\label{prop:3.1}
$d$Laplace$\xi$
\begin{equation}
    \varepsilon_{\PE} \leq I(Y; P \mid X) \cdot \exp\!\left(-\frac{d}{d_0}\right)
\end{equation}
$d_0 = O(L/\xi)$
\end{proposition}

\heuristic{} \textbf{}1.2.2
PinskerFourier
Cauchy$\propto \exp(-d/d_0)$——
$P_{Y|X,P}$$O(1/d)$

% ---------------------------------------------------------------------------
\subsection{Theorem 2——}

\begin{theorem}[Theorem $2'$——\Situs{}]
\label{thm:2'}
$X$$\delta > 0$$\varepsilon_{\PE}$\Situs{}
\Situs{}SCX$F_1$
\begin{equation}
    \boxed{
    F_{1,\text{SCX}}^{\Situs} \leq F_{1,\text{base}}
    + C_F \cdot \sqrt{\frac{\delta + \frac{2\varepsilon_{\PE}}{C_F^2}}{2}}
    }
\end{equation}
$C_F > 0$$F_1$
\end{theorem}

\begin{proof}
\textbf{1Fano}
$P_e \geq \frac{H(Y \mid X, \PE(P)) - 1}{\log |\Y|}$

\textbf{2}
\begin{align}
    H(Y \mid X, \PE(P))
    &= H(Y \mid X) - I(Y; \PE(P) \mid X) \nonumber\\
    &= H(Y \mid X) - [I(Y; P \mid X) - \varepsilon_{\PE}] \nonumber\\
    &= H(Y \mid X) - I(Y; P \mid X) + \varepsilon_{\PE}
\end{align}

\textbf{3$\delta$}$\delta_{\Situs} = \delta + 2\varepsilon_{\PE}/C_F^2$
\Situs{}

\textbf{4SCX$F_1$}
$F_{1,\text{SCX}}^{\Situs} \leq F_{1,\text{base}} + C_F \cdot \sqrt{\delta_{\Situs}/2}$

\textbf{5}$\varepsilon_{\PE} \to 0$
Theorem~2$\hfill \square$
\end{proof}

\rigorous{} \textbf{}
$C_F$$F_1$$F_1$
$F_1 = 2\text{TP}/(2\text{TP} + \text{FP} + \text{FN})$


\textbf{}
\begin{itemize}
    \item $\varepsilon_{\PE} \to 0$\textbf{Theorem 2}——
          $F_{1,\text{SCX}}^{\Situs} \leq F_{1,\text{base}} + C_F\sqrt{\delta/2}$
    \item $\varepsilon_{\PE} \to I(Y; P \mid X)$
          ——\textbf{}
\end{itemize}

% ---------------------------------------------------------------------------
\subsection{$\varepsilon_{\PE}$}

\begin{theorem}[$\varepsilon_{\PE}$——KSG]
\label{thm:3.3.1}
$\varepsilon_{\PE}$\textbf{}
\begin{equation}
    \boxed{
    \hat{\varepsilon}_{\PE}^{(n)} = \hat{I}^{(n)}(Y; P \mid X)
                                 - \hat{I}^{(n)}(Y; \PE(P) \mid X)
    }
\end{equation}
$\hat{I}^{(n)}(\cdot; \cdot \mid \cdot)$Kraskov-St\"{o}gbauer-Grassberger (KSG)
$n \to \infty$$\hat{\varepsilon}_{\PE}^{(n)} \xrightarrow{p} \varepsilon_{\PE}$
\end{theorem}

\textbf{}
\begin{enumerate}[label=(\arabic*)]
    \item $n$$(x_i, p_i, y_i)$
    \item $\PE(p_i)$
    \item $\hat{I}^{(n)}(Y; (P, X))$$(\Ppos \times \X) \times \Y$$k$-NN
    \item $\hat{I}^{(n)}(Y; X)$$\X \times \Y$$k$-NN
    \item $\hat{I}^{(n)}(Y; P \mid X) = \hat{I}^{(n)}(Y; (P, X)) - \hat{I}^{(n)}(Y; X)$
    \item $\hat{I}^{(n)}(Y; \PE(P) \mid X)$
    \item $\hat{\varepsilon}_{\PE}^{(n)} = \hat{I}^{(n)}(Y; P \mid X) - \hat{I}^{(n)}(Y; \PE(P) \mid X)$
\end{enumerate}

\textbf{}KSG$O(n^{-1/(d_{\text{eff}}+2)})$$d_{\text{eff}}$
\heuristic{} $X$——$\varepsilon_{\PE}$
\textbf{}

\begin{proposition}[$\varepsilon_{\PE}$]
\label{prop:3.3.1}
$f_{\text{full}}$$(X, P)$$f_{\text{enc}}$$(X, \PE(P))$
\textbf{}$\varepsilon_{\PE}$
\begin{equation}
    \boxed{
    \hat{\varepsilon}_{\PE}^{\text{pred}} = \frac{1}{n_{\text{test}}}
    \sum_{i=1}^{n_{\text{test}}}
    \left[\log\frac{1}{p_{\text{enc}}(y_i \mid x_i, \PE(p_i))}
        - \log\frac{1}{p_{\text{full}}(y_i \mid x_i, p_i)}\right]
    }
\end{equation}
$n_{\text{test}} \to \infty$$\hat{\varepsilon}_{\PE}^{\text{pred}} \to \varepsilon_{\PE}$
\end{proposition}

\heuristic{} \textbf{}\textbf{}
$\varepsilon_{\PE} \approx 0$——\Situs{}Theorem~2
KSG
——$\varepsilon_{\PE}$

% ===========================================================================
% SECTION 5: THEOREM 3 CORRECTION
% ===========================================================================
\section{Theorem $3'$}

% ---------------------------------------------------------------------------
\subsection{Theorem 3}

\begin{theorem}[Theorem 3——]
\label{thm:3_original}
$W_A, W_B$
\begin{enumerate}[label=(\arabic*)]
    \item $P_{W_A}(X, Y) = P_{W_B}(X, Y)$
    \item $W_A$$(x, y)$\textbf{}label
    \item $W_B$$(x, y)$\textbf{}labellabel
\end{enumerate}
i.i.d.$(X_i, Y_i)_{i=1}^n$$W_A$$W_B$
$d: (\X \times \Y)^n \to \{W_A, W_B\}$
\begin{equation}
    \max_{W \in \{W_A, W_B\}} P_W(d(\text{data}) \neq W) \geq \frac{1}{2}
\end{equation}
\end{theorem}

3-``''
``''``''
——

% ---------------------------------------------------------------------------
\subsection{PE}

\Situs{}$(X, P, Y)$$P$
\textbf{}PE

\begin{proposition}[PETheorem 3]
\label{prop:4.1}
$\PE: \Ppos \to \R^{d_{\pe}}$\textbf{}

\begin{equation}
    P_{W_A}(X, P, \PE(P), Y) = P_{W_B}(X, P, \PE(P), Y)
\end{equation}
Theorem~3\textbf{}
\end{proposition}

\begin{proof}
$\PE$$(X, P, \PE(P), Y)$$(X, P, Y)$
$P$
$P_{W_A}(P \mid X) = P_{W_B}(P \mid X)$
$P_{W_A}(X, Y) = P_{W_B}(X, Y)$Theorem~3
\hfill $\square$
\end{proof}

\rigorous{} \textbf{}``''


% ---------------------------------------------------------------------------
\subsection{PE——}

\textbf{}Theorem~3$P_{W_A}(X, Y) = P_{W_B}(X, Y)$
\textbf{}
\begin{itemize}
    \item $W_A$$Y_{\text{true}}$$X$$Y_{\text{obs}}$$\eta$
    \item $W_B$$Y_{\text{obs}} = Y_{\text{true}}$$P(Y_{\text{true}} \mid X)$
          
\end{itemize}

PE$\theta$$Y$$Y_{\text{true}}$
——

\begin{theorem}[PETheorem 3]
\label{thm:4.1}
$\mathcal{L}_{\text{total}}(\theta) = \mathcal{L}_{\text{sup}}(\theta)
+ \lambda \cdot \mathcal{L}_{\text{aux}}(\theta)$
\begin{itemize}
    \item $\mathcal{L}_{\text{sup}}$$(X, P, Y)$
    \item $\mathcal{L}_{\text{aux}}$$Z$
\end{itemize}
$\mathcal{L}_{\text{aux}}$$W_A$$W_B$
\begin{equation}
    \argmin_\theta \mathcal{L}_{\text{aux}}^{W_A}(\theta)
    \neq \argmin_\theta \mathcal{L}_{\text{aux}}^{W_B}(\theta)
\end{equation}
PE$\hat{\theta}_{W_A} \neq \hat{\theta}_{W_B}$1
$n \to \infty$\textbf{}
\end{theorem}

\textbf{}
\begin{enumerate}[label=(\roman*)]
    \item \textbf{}PE$Y_{\text{true}}$$Z$
          $Z$
    \item \textbf{}$W_A$
          $W_B$
    \item \textbf{}PE
\end{enumerate}

% ---------------------------------------------------------------------------
\subsection{Theorem $3'$——}

\begin{theorem}[Theorem $3'$——]
\label{thm:3'}
\begin{quote}
\textbf{ (A):} PE$Y$\\
\textbf{ (B):} $P$$X$
$P \indep W \mid X$\\
\textbf{:} Theorem~3$d$
$\max_{W \in \{W_A, W_B\}} P_W(d \neq W) \geq 1/2$\\
\textbf{ (A) :} \textbf{Open Problem}
PE
——\S5.3.2
\end{quote}
\end{theorem}

\rigorous{} \textbf{(A)(B)
(A)``''
——``''\S5.3.2Theorem~3
PE
}

% ---------------------------------------------------------------------------
\subsection{}
\label{sec:4.3}

Theorem~$3'$

\subsubsection{}

$X \in \{-1, +1\}$$P \in \{0, 1\}$$Y \in \{-1, +1\}$

\textbf{$W_A$}
\begin{itemize}
    \item $P(X=+1) = P(X=-1) = 1/2$
    \item $Y_{\text{true}} = X$
    \item $\eta = 0.2$$Y_{\text{obs}} = Y_{\text{true}}$$0.8$
          $Y_{\text{obs}} = -Y_{\text{true}}$$0.2$
    \item $P = X$
\end{itemize}

\textbf{$W_B$}
\begin{itemize}
    \item $P(X=+1) = P(X=-1) = 1/2$
    \item $P(Y=+1 \mid X=+1) = 0.8$, $P(Y=-1 \mid X=+1) = 0.2$
    \item $P(Y=+1 \mid X=-1) = 0.2$, $P(Y=-1 \mid X=-1) = 0.8$
    \item $Y_{\text{obs}} = Y_{\text{true}}$
    \item $P = X$$W_A$
\end{itemize}

\textbf{}$P_{W}(X, Y)$

\subsubsection{PE}

PE$\PE_\theta(p) = \theta \cdot p$$\theta \in \R$
$\mathcal{L}_{\text{aux}}(\theta) = \E[(\PE_\theta(p) - \E[Y \mid X=x])^2]$

\textbf{}$\E_{W_A}[Y \mid X] = \E_{W_B}[Y \mid X] = 0.6 \cdot \text{sign}(x)$
——\textbf{}\textbf{PE}
Theorem~3

\textbf{Theorem~3}``''
Theorem~3'\textbf{}PE\textbf{}


\openproblem{} \textbf{3``''}
\textbf{}$Z$——$W_A$$|Z_{\text{pred}} - Z_{\text{DFT}}|$
$W_B$——PE
\textbf{}\textbf{}
$Z$$\mathcal{L}_{\text{aux}}$


\subsubsection{}

\begin{corollary}[]
\label{cor:4.4.1}
\begin{enumerate}[label=(\arabic*)]
    \item \textbf{PEPE}Theorem~3
          ——\textbf{}
    \item \textbf{PE}$Y$Theorem~3
    \item \textbf{PE}Theorem~3\textbf{}
          \textbf{Open Problem}——
          ``''
    \item \textbf{}
          ``''
\end{enumerate}
\end{corollary}

\heuristic{} (3)``''
PECC

% ===========================================================================
% SECTION 6: DISCUSSION
% ===========================================================================
\section{}

% ---------------------------------------------------------------------------
\subsection{}

\subsubsection{RoPE/}

RoPE\cite{su2024roformer}\cite{vaswani2017attention}
\textbf{}$\omega_j = 10000^{-2j/d}$
token

\Situs{}1.2.1$\omega_j \propto \tan(\pi(2j+1)/2d)$

~\ref{tab:comparison}

\begin{table}[htbp]
\centering
\caption{PEPE}
\label{tab:comparison}
\begin{tabular}{@{}p{3.5cm}ll@{}}
\toprule
\textbf{} & \textbf{RoPE/} & \textbf{\Situs{}} \\
\midrule
 &  &  \\
 &  & Laplace $\exp(-|\Delta|/\xi)$ \\
 &  & Bochner +  \\
Lipschitz &  & 1.2.3/1.3.1 \\
 &  & $\xi$ \\
 & RoPE: 2D & \Situs{} 3D: \textbf{} \\
\bottomrule
\end{tabular}
\end{table}

\subsubsection{——}

\Situs{}3D1.3.1\textbf{}
$R_{\text{phys}} \in SO(3)$$\mathbf{p} \to R_{\text{phys}}\mathbf{p}$
$\mathbf{p}$$T$$T \cdot \PE(\mathbf{p}) = \PE(R_{\text{phys}}\mathbf{p})$
——$x$$y$
``''

/
slab

\textbf{}
$\PE(\mathbf{p}) = [Y_{\ell m}(\theta, \phi)]$$SO(3)$
E(3)Tensor Field Networks\cite{thomas2018tfn}
SE(3)-Transformers\cite{fuchs2020se3}$(L+1)^2$
$L=3$16 vs 6Lipschitz

\openproblem{} \textbf{1}$\R^3 \to SO(d)$
1.3.1$SO(3)$——
$\R^3 \to SO(d)$Abel$SO(3)$Abel

% ---------------------------------------------------------------------------
\subsection{}

\textbf{}——/
\Situs{}``$i$$p_i$''$i$
$h_i$——


\begin{itemize}
    \item \textbf{}DFT
    \item \textbf{}
          
\end{itemize}

\textbf{}
——

\openproblem{} \textbf{2}\Situs{}
$h_i = \phi(s_i) + \PE(p_i)$

% ---------------------------------------------------------------------------
\subsection{\Situs{}——}

\Situs{}\textbf{}

\begin{enumerate}[label=(F\arabic*), leftmargin=*]
    \item \textbf{$I(Y; P \mid S) = 0$}
          2.2.1$\delta_s^{\PE} \to 0$\Situs{}$d$
    \item \textbf{}$\Ppos$
          →artifact
    \item \textbf{/}
          \Situs{}→
    \item \textbf{}GNN
          pLM→\Situs{}\S2.2 of
          situs\_physical\_validation.md
    \item \textbf{+}$I(Y; P \mid S)$→
          $\delta_s^{\PE} > 0$$\delta_s^{\PE} < 0$
    \item \textbf{artifact}DFT→
          $N$artifact\S1.3 of
          situs\_physical\_validation.md
\end{enumerate}

% ---------------------------------------------------------------------------
\subsection{}

~\ref{tab:honesty}/

\begin{table}[htbp]
\centering
\caption{}
\label{tab:honesty}
\begin{tabular}{@{}lll@{}}
\toprule
\textbf{/} & \textbf{} & \textbf{} \\
\midrule
1.2.1 & Bochner+ & \rigorous  \\
1.2.2 &  & \rigorous  \\
1.2.3 & Lipschitz & \rigorous  \\
1.3.1 & Lipschitz & \rigorous  \\
1.4.1 & Lipschitz & \rigorous  \\
\midrule
2.1 &  & \rigorous  \\
2.2.1 & $\delta_s^{\PE}>0$ & \rigorous 1--3; \heuristic 4 \\
2.3.1 & DPI & \rigorous $M$ \\
2.4.1 & Fano & \heuristic  \\
2.5.1 &  & \rigorous  \\
\midrule
3.1 & $\varepsilon_{\PE}$ & \rigorous  \\
3.1 & - & \heuristic  \\
Theorem $2'$ &  & \rigorous  \\
3.3.1 & KSG & \rigorous ; \heuristic  \\
3.3.1 &  & \heuristic  \\
\midrule
4.1 & PETheorem 3 & \rigorous  \\
4.1 & PE & \openproblem  \\
Theorem $3'$ &  & \rigorous ; \openproblem  \\
\bottomrule
\end{tabular}
\end{table}

\textbf{}
\begin{enumerate}[label=(\arabic*)]
    \item \textbf{}\S6.1.23D$SO(3)$
    \item \textbf{}2.2.14
    \item \textbf{2.4.1Fano}``''Fano——
\end{enumerate}

% ---------------------------------------------------------------------------
\subsection{Open Problem}

\begin{enumerate}[label=(\arabic*), leftmargin=*]

    \item \textbf{Chernoff}Multi-Head Spring
          Theorem~1Chernoff bound$\beta$-mixing
          \cite{scx2026cc_audit}

    \item \textbf{$SO(3)$}$\R^3 \to SO(d)$
          $SO(3)$
          

    \item \textbf{\Situs{}}
          

    \item \textbf{$\varepsilon_{\PE}$}KSG
          bootstrapsubsampling
          $\varepsilon_{\PE} = 0$``\Situs{}''

    \item \textbf{}1.2.1$\xi$
          $\xi$-
          

    \item \textbf{Theorem~$3'$}
          PE
          

    \item \textbf{\Situs{}}GNNpLM
          \Situs{}``''
          \Situs{}\emph{}

\end{enumerate}

% ===========================================================================
% ACKNOWLEDGMENTS
% ===========================================================================
\section*{}

SCXCC\cite{scx2026cc_audit}
Theorem~1--4
$\delta_s^{\PE}$CC2.1sign error
$\varepsilon_{\PE}$3.1CC
\texttt{ppe\_rigorous\_derivation.md}

% ===========================================================================
% BIBLIOGRAPHY
% ===========================================================================
\begin{thebibliography}{99}

\bibitem{vaswani2017attention}
A.~Vaswani, N.~Shazeer, N.~Parmar, J.~Uszkoreit, L.~Jones, A.~N.~Gomez,
L.~Kaiser, and I.~Polosukhin.
\newblock Attention is all you need.
\newblock In \emph{Advances in Neural Information Processing Systems (NeurIPS)},
  2017.

\bibitem{su2024roformer}
J.~Su, M.~Ahmed, Y.~Lu, S.~Pan, W.~Bo, and Y.~Liu.
\newblock Ro{Former}: Enhanced transformer with rotary position embedding.
\newblock \emph{Neurocomputing}, 568:127063, 2024.

\bibitem{devlin2019bert}
J.~Devlin, M.-W.~Chang, K.~Lee, and K.~Toutanova.
\newblock {BERT}: Pre-training of deep bidirectional transformers for language
  understanding.
\newblock In \emph{Proceedings of NAACL-HLT}, 2019.

\bibitem{scx2026theorems}
SCX Project.
\newblock {Theorem 1--4}: Multi-expert consistency, weak feature limits,
  unidentifiability, and minimax optimality.
\newblock Technical report, \texttt{theory/theorems/}, 2026.

\bibitem{scx2026cc_audit}
SCX Project.
\newblock Multi-head spring and positional encoding analysis: {CC} audit report.
\newblock Technical report,
  \texttt{theory/self\_evolution/multi\_head\_spring\_and\_positional\_encoding\_analysis.md},
  2026.

\bibitem{scx2026ppe_derivation}
SCX Project.
\newblock Physical positional encoding ({PPE}) rigorous derivation in the {SCX}
  framework.
\newblock Technical report,
  \texttt{theory/self\_evolution/ppe\_rigorous\_derivation.md}, 2026.

\bibitem{scx2026situs_validation}
SCX Project.
\newblock {Situs} ({Physical Positional Encoding}) physical validation and
  applicability analysis.
\newblock Technical report, \texttt{theory/self\_evolution/situs\_physical\_validation.md},
  2026.

\bibitem{thomas2018tfn}
N.~Thomas, T.~Smidt, S.~Kearnes, L.~Yang, L.~Li, K.~Kohlhoff, and P.~Riley.
\newblock Tensor field networks: Rotation- and translation-equivariant neural
  networks for {3D} point clouds.
\newblock In \emph{NeurIPS}, 2018.

\bibitem{fuchs2020se3}
F.~B.~Fuchs, D.~E.~Worrall, V.~Fischer, and M.~Welling.
\newblock {SE(3)-Transformers}: {3D} roto-translation equivariant attention
  networks.
\newblock In \emph{NeurIPS}, 2020.

\bibitem{wainwright2019}
M.~J.~Wainwright.
\newblock \emph{High-Dimensional Statistics: A Non-Asymptotic Viewpoint}.
\newblock Cambridge University Press, 2019.

\bibitem{boucheron2013}
S.~Boucheron, G.~Lugosi, and P.~Massart.
\newblock \emph{Concentration Inequalities: A Nonasymptotic Theory of
  Independence}.
\newblock Oxford University Press, 2013.

\bibitem{freysoldt2014}
C.~Freysoldt, B.~Grabowski, T.~Hickel, J.~Neugebauer, G.~Kresse,
A.~Janotti, and C.~G.~Van de Walle.
\newblock First-principles calculations for point defects in solids.
\newblock \emph{Reviews of Modern Physics}, 86(1):253--305, 2014.

\bibitem{makov1995}
G.~Makov and M.~C.~Payne.
\newblock Periodic boundary conditions in ab initio calculations.
\newblock \emph{Physical Review B}, 51(7):4014--4022, 1995.

\bibitem{pauling1951}
L.~Pauling, R.~B.~Corey, and H.~R.~Branson.
\newblock The structure of proteins: Two hydrogen-bonded helical configurations
  of the polypeptide chain.
\newblock \emph{Proceedings of the National Academy of Sciences},
  37(4):205--211, 1951.

\bibitem{dunker2005}
A.~K.~Dunker, M.~S.~Cortese, P.~Romero, L.~M.~Iakoucheva, and V.~N.~Uversky.
\newblock Flexible nets: The roles of intrinsic disorder in protein interaction
  networks.
\newblock \emph{The FEBS Journal}, 272(20):5129--5148, 2005.

\bibitem{cover2006}
T.~M.~Cover and J.~A.~Thomas.
\newblock \emph{Elements of Information Theory}, 2nd edition.
\newblock Wiley, 2006.

\bibitem{ksg2004}
A.~Kraskov, H.~St\"{o}gbauer, and P.~Grassberger.
\newblock Estimating mutual information.
\newblock \emph{Physical Review E}, 69(6):066138, 2004.

\end{thebibliography}

% ===========================================================================
% APPENDIX (optional — reserved for extended proofs)
% ===========================================================================
% \appendix
% \section{}
% ()

\end{document}
